\begin{auxlemma}[Evaluation of the cell norm at a positive Gram point]
\label{lem:unit-cell-evaluation}
Let \(J_\mu\) be the amplitude in Lemma~\ref{lem:minor-orbital-shift},
\(\mu=\lambda+\nu\), \(\|\nu\|_1\le C_nkm^3\).
Use the actual ordered words, compact types and root measures at a
marked real place. Write \(w(\gamma)\) for the orbital Jacobian of
the block containing \(\gamma\), and \(G=G_{Z,v}>0\).
For all sufficiently late marked rows, uniformly for every nonzero
pair of periods in this window,
\begin{equation}\label{eq:unweighted-unit-cell-evaluation}
 w(\gamma)\,\frac{\|J_\mu(\gamma)\|_{\rm HS}^2}{M_\mu}
       \le b^{C_bm^4}\max(1,G^{-1})^{C_bm^2},
 \qquad M_\mu=\|J_\mu\|_{\rm cell}^2 .
\end{equation}
The constant \(C_b\) depends on the fixed weight, character and
measure bounds. The threshold may also depend on \(C_n\).
\end{auxlemma}

\begin{proof}
\proofstep{Point evaluation on a sequence of blocks.}
Use the single compact-finite coefficient representing \(J_\mu\)
in Proposition~\ref{prop:local-operator-normalization}. For any such
coefficient \(F\), denote its transformed squared norm at the point
and its integrated squared norm on the big cell by
\[
 E_F(\gamma)=w(\gamma)\|J_F(\gamma)\|_{\rm HS}^2,\qquad
 {\cal M}(F)=\int_{\text{big cell}} E_F(c)\,dc .
\]
Put \(p=r+1\), \(p_0=2p\), \(N=3p\). In a cut between differing
compact coordinates divide a diagonal word energy by its positive
word weights as in Proposition~\ref{prop:real-node-positivity}.
For the point estimate one needs the three-derivative profile
\[
             {\cal P}(\xi)=\min\{1,(\Lambda/|\xi|)^3\},
             \qquad {\cal P}(0)=1,\qquad \Lambda\le b^{C m}.
\]

A cut of length \(30p^2+6\) admits \(3p(2r+2)=6p^2\) nodes
five units apart and at least two units from the endpoints, since their total required length is \(5(6p^2-1)+4\). In each interpolation take fresh
nodes for the degree (at most \(2r+1\)) of the new letter, using
Proposition~\ref{prop:letter-polynomial-degree}. Use the positive
normalizer \(\prod_i(|Y+\lambda_i|+c_0)\). With \(n_*=30p^2+6\)
and \(A_*=4p+n_*\), one may take
\[
 e=\|\nu\|_1+p(A_*+n_*+N+2),\qquad c_0=2+2e .
\]

Multiplication by the determinant costs \(pA_*\) unit shifts, the
compact contractions cost at most \(pn_*\), and \(N\) radial letters
use at most \(pN\) further shifts. Thus the normalizer includes every
ambient eigenvalue used below. Writing \(h_*\) for the least
unshifted row gap, Lemma~\ref{lem:marked-paired-moment-range} gives
\[
 h_* \ge \min(4\Delta,\eta_{\rm size}/m,b_1)>nm^{10},
 \qquad e\le C_nkm^3+O(m^3),\quad 200 c_0<h_*
\]
on the late family. In particular all shift and degree budgets in
Definition~\ref{hyp:paired-frame-moment-range} allow these endpoints.

For each of the two column polynomials, pair the \(k-1\) interlaced
row roots with the adjacent column roots \emph{away from} the cut.
Each such numerator distance is at most its paired denominator
distance. Omit the two coordinates adjoining the cut if the compact
type has been moved there. Each exceptional numerator is bounded
by \(b^{O(1)}\), and the interior nodes stay a fixed distance from
the column poles and their unit translates.

Padding the row roots by \(c_0\) has a bound even at the actual
endpoint radius \(pq\). Remove a root nearest \(Y\), and bound its
padded factor directly by \(b^{O(1)}\). The \(j\)-th remaining
root on either side has distance at least \(jh_*/2\). Thus the
other ratios have product at most
\[
 \exp\!\left(\frac{4c_0}{h_*}\sum_{j=1}^n\frac1j\right)
        \le (\mathrm{e}n)^{4c_0/h_*}\le\mathrm{e}n .
\]
The other \(m\) ambient factors are each at most \(b^{O(1)}\).
These estimates hold also at the adjacent unit translates of the
nodes. Previous letters change the rational factors there by bounded
ratios for each of the at most \(N\) letters, since distinct nodes
are five units apart. Interpolating a leading coefficient on
integer-spaced nodes has total absolute cost at most \(2^{2r+1}\),
by the product formula for the interpolation weights. This gives
the asserted \(\Lambda\), after incorporating the character and line
norms of Definition~\ref{hyp:paired-frame-moment-range}.

Apply the signed derivative step of
Proposition~\ref{prop:positive-block-fourier} three times, as in
the point recursion following Lemma~\ref{lem:empty-word-exhaustion}. For a size-one block, Fourier inversion of its diagonal
costs at most
\[
              \int_{\C}{\cal P}(\xi)\,d\xi=3\pi\Lambda^2 .
\]
For a larger block choose the coset of the square lattice of side
\(\Lambda^{-1}\) through the prescribed point. Keep that diagonal
in the nonnegative period sum of Lemma~\ref{lem:germ-period}.
The comparison for a block of one smaller size uses the
inverse-coordinate frequencies in Lemma~\ref{lem:gram-limit}.
The \(j\)-th annulus of the dual lattice contains \(O(j+1)\)
frequencies. For a block of size \(s\) this gives
\[
 \sum_{\xi\ {\rm dual}}{\cal P}(\xi)
       \ll 1+\sum_{j\ge1}(j+1)j^{-3}\ll1,\qquad
 E_{h\|(s,t)}\le(C\Lambda^2)^s\max_{|w'|-|w|\le3s}E_{h,w'} .
\]
The power of \(\Lambda^2\) at each unit of size is the reciprocal
covolume of the translation lattice.

Apply this argument first away from \(G=0\) with a smooth cutoff.
At total size \(p\) the cost is \((C\Lambda^2)^p\). At the empty
prefix the limiting normalized energy of every resulting word is
at most \({\cal M}(F)\) when \(F\) is the coefficient multiplied
by \(G^{p_0}\) and the compact contractions below. Indeed the
commuting Casimirs can be moved to this coefficient in the big-cell
pairing by Lemma~\ref{lem:radial-cutoff-transfer}. Tensoring with the
matrix polynomials in \(F\) gives a finite sum of infinitesimal
characters. Different characters are orthogonal under the invariant
pairing, and the positive normalizer bounds each eigenvalue at
every letter. Hence its normalized product has modulus at most one.

The extra matrix polynomial and the cutoff satisfy
\[
 \deg(\text{multiplier})\le 2pp_0+pn_0,\qquad n_0\le30p^2+6,\qquad
                        N=3p\le4p=2p_0 .
\]
The comparison following Lemma~\ref{lem:empty-word-exhaustion}
gives the integrable majorant for both compact types. For a fixed positive
point the cutoff eventually equals one. Removing it gives
\[
                    E_F(\gamma)\le(C\Lambda^2)^p{\cal M}(F).
\]

\proofstep{Obtaining a cut for both compact types.}
For different marked types a differing coordinate gives a cut of
length at least \(30p^2+6\) by
Definition~\ref{hyp:critical-level-weights} and the marked spacings.
Use \(F=G^{p_0}f\) with \(f\) representing \(J_\mu\), and put
\(n_0=0\). When the types agree take \(n_0=30p^2+6\) and choose
a gap in the compact row. The shifts producing the two long cuts are
\[
        u_i-u_{i+1}\ge2n_0+1,\qquad
    u^+=u+n_0\mathbf1_{\{i+1\}},\quad u^-=u-n_0\mathbf1_{\{i\}} .
\]

For clarity, let \(L_1>\dots>L_k\) be the strictly decreasing integer
coordinates of that compact row after subtracting its common scalar.
On restriction to the stabilizer of a unit vector the type labels
\(M\) satisfy \(L_{j+1}<M_j\le L_j\). The squared projections when
one upper coordinate \(L_a\) is raised or lowered are respectively
\[
 \frac{\prod_{j<k}|M_j-1-L_a|}{\prod_{j\ne a}|L_j-L_a|},
 \qquad
 \frac{\prod_{j<k}|M_j-L_a|}{\prod_{j\ne a}|L_j-L_a|}.
\]
These are the residues of the compact branching polynomials in
Lemma~\ref{lem:residual-single-cofactors}: taking the \(a\)-th
spectral projection divides the characteristic polynomial of the
compressed operator by the derivative of the full polynomial at
\(L_a\). The raising polynomial has roots \(M_j-1\) and the lowering
polynomial roots \(M_j\), from the unit shift in the dual action.

For \(M_i\ge L_{i+1}+n_0+1\) the entire raising chain is allowed.
Otherwise \(M_i\le L_i-n_0\) and the lowering chain is allowed.
In any allowed step write \(b_j=\max(1,|L_j-L_a|)\).
The \(j\)-th numerator factor is at least \(b_{j+1}\) for \(j<a\)
and \(b_j\) for \(j\ge a\), by interlacing and integrality.
Multiplying the lower bounds gives the cancellation
\[
 \frac{\prod_{j<a}b_{j+1}\prod_{a\le j<k}b_j}
      {\prod_{j\ne a}b_j}
       =\frac{1}{b_1b_k},\qquad b_a=1 .
\]
If \(B_1=b^{C'}\) bounds the compact entries
and the shifted width is at most \(3B_1\), each squared projection
on the allowed chain is at least \((3B_1)^{-2}\).
Both chains have operator norm at most one at a unit vector.

To obtain the contraction vector at \(\gamma\), take the last row of the
adjugate of the cross-pairing matrix. Combine the left frame with
its conjugate coefficients, transport by \(\gamma^{-1}\), then
project onto the right positive middle space perpendicular to \(z'\).
The adjugate identity annihilates its pairings with the first \(r\)
right columns. Its pairing with \(z'\) has modulus \(|\Delta_p|\).
Before subtracting its \(z'\)-component only the left distinguished
line contributes to the norm. Its squared positive length is thus
\[
               |\Delta_p|^2/T'-T|\Delta_r|^2=G/T' .
\]

Let \(A_\pm\) be the two contraction chains evaluated at this vector,
oriented from the shifted compact type to the original one. Sum the
squared lower bound over the orthogonal \(M\)-components and use
homogeneity of degree \(n_0\). For any \(y\) in the original type,
\[
 \|A_\pm\|^2\le(G/T')^{n_0},\qquad
 \sum_\pm\|A_\pm^*y\|^2\ge (3B_1)^{-2n_0}(G/T')^{n_0}\|y\|^2 .
\]
The contraction is invariant under both unipotents and the left
stabilizer. Right compact translation changes \(A_\pm\) by its
two compact actions, which cancel those in the orbital integral.
Thus \(J_{F_\pm}=G^{p_0}J_f A_\pm\) for \(F_\pm=G^{p_0}f A_\pm\).
The same point, stabilizer Haar and quotient measure occur in all
three integrals. Taking adjoints and summing over a target
orthonormal basis gives
\[
 E_f(\gamma)\le(3B_1)^{2n_0}(T')^{n_0}G^{-A}
                 \sum_\pm E_{F_\pm}(\gamma),\qquad A=2p_0+n_0.
\]
For different types the same estimate has just one term and \(A=2p_0\).

\proofstep{The determinant moment in the integrated norm.}
On the big cell \(0<G\le|\Delta_p|^2\). Hence
\[
 {\cal M}(F_\pm)\le (T')^{-n_0}
       \int_{\text{big cell}}|\Delta_p|^{2A}E_f(c)\,dc .
\]
Use \eqref{eq:minor-shift-unit} with \(l=p\), now \(A\) times,
in the Hilbert space of square-integrable operator-valued sections.
Each nonzero endpoint has norm one. Minkowski's inequality gives
\[
 \|\Delta_p^A U_\mu\|_{\rm cell}\le B^A,\qquad
 \int_{\text{big cell}}|\Delta_p|^{2A}E_f(c)\,dc\le B^{2A}M_\mu .
\]
Here \(A\le40p^2\), so the extra radius \(Ap\le40p^3\) fits
the window of Lemma~\ref{lem:minor-orbital-shift} after increasing
\(C_n\) by one. This use of its finite expansion precedes any
bound on a point value. The factors \((T')^{n_0}\) cancel.
The two remaining factors satisfy
\[
    2(C\Lambda^2)^p(3B_1)^{2n_0}B^{2A}\le b^{O(m^4)},\qquad
                   G^{-A}\le\max(1,G^{-1})^{10m^2}.
\]

Finally, transporting the character to unit size by the flag torus
scales its \(j\)-th column by \(\alpha^j\) for \(0\le j\le r\).
The line and Gram scalings are
\[
 (T,T')\longmapsto\alpha^{2r}(T,T'),\qquad
    G=|\Delta_p|^2-TT'|\Delta_r|^2\longmapsto\alpha^{2rp}G .
\]

For a prefix of length \(a\) the minor scales by
\(\alpha^{a(a-1)}\). In each ordinary block the \(2k\) part of the
density exponent multiplies the ratio of the two block-end minors.
Across the blocks these terms telescope. If the last block is a
singleton the identity \(G=|\Delta_r|^2(|t_0|^2-TT')\) gives the
same total exponent \(2krp\) as on the big cell. The common
\(k\)-dependent dilation cancels in the ratio of the two energies.

In big-cell pivot coordinates the \(j\)-th complex coordinate
scales by \(\alpha^{2j}\), so its section Jacobian is
\[
             J_{\rm sec}=\prod_{j=0}^r\alpha^{4j}=\alpha^{2rp}.
\]

In a block of size \(s\) the within-block root coordinates are
indexed by pairs \(a<b\). Write \(D_{\rm rt}=\sum(b-a)\) over these pairs and
\(c=\sum_{\rm blocks}\binom{s}{2}\). On a block starting after
\(j\) columns the remaining density exponent, relative to splitting
the block into singletons, is
\[
 4\sum_{i=0}^{s-1}(j+i)(2i+1-s)
       =4\sum_{0\le a<b<s}(b-a).
\]

The root \(ab\) has complex-coordinate Jacobian
\((\alpha^b/\alpha^a)^2\) and contributes \(\alpha^{-1}\)
in the unit-root normalization. Thus the square of their
combined orbital change has exponent \(4D_{\rm rt}-2c\).
Canceling against the density and section changes leaves the
factor \(J_{\rm sec}\alpha^{-2c}\), or its reciprocal on reversal. Both therefore satisfy
\[
 \max\{J_{\rm sec}\alpha^{-2c},(J_{\rm sec}\alpha^{-2c})^{-1}\}
   \le\max(\alpha,\alpha^{-1})^{2rp+2c},\qquad 2c\le rp .
\]

Multiplication by \(G^{-A}\) costs at most a further \(2rpA\)
powers of this scale. The independent root Haar factor \(\theta'\)
in \eqref{eq:absolute-root-scales} satisfies
\(b^{-|C_1|m^2}\le\theta'\le b^{|C_1|m^2}\), so its inverse square
costs at most \(b^{2|C_1|m^2}\). Since
\(\max(\alpha,\alpha^{-1})\le b^{O(1)}\), all these powers are
absorbed in \(b^{C_bm^4}\).
\end{proof}
